\section{Konsep Dasar Serangan}
\label{sec:konsep-serangan}

Tujuan utama kriptografi adalah menjaga kerahasiaan pesan atau kunci dari pihak
ketiga, yang dalam konteks ini disebut sebagai musuh, lawan, penyadap, atau
kriptanalis. Bab sebelumnya memperkenalkan keempat layanan keamanan dan tiga
keluarga algoritma; bab ini membahas sisi seberangnya, yaitu bagaimana usaha-usaha
tersebut dapat digagalkan. Sebagaimana disinggung pada Bagian~\ref{sec:kriptanalisis},
kriptografi dan kriptanalisis tumbuh bersama-sama, sehingga pemahaman yang utuh
menuntut keduanya dipelajari berdampingan.

Sebuah \textbf{serangan} (\textit{attack}) diartikan sebagai setiap usaha
(\textit{attempt}) atau percobaan yang dilakukan pihak lawan untuk menemukan kunci
atau menemukan plainteks dari cipherteksnya. Definisi ini sengaja luas, mencakup
percobaan yang berhasil maupun yang gagal. Yang dinilai pada akhirnya bukan apakah
sebuah serangan pernah dicoba, melainkan berapa besar sumber daya --- waktu,
komputasi, memori, atau biaya --- yang dibutuhkan untuk membuatnya berhasil
\citep{katz2020}.

Dalam menganalisis serangan, kita menggunakan \textbf{Prinsip Kerckhoffs} yang
sudah diperkenalkan pada Bagian~\ref{sec:terminologi}:
\begin{quote}
    \textit{Semua algoritma kriptografi harus publik; hanya kunci yang rahasia.}
\end{quote}
Oleh karena itu, keamanan sistem kriptografi tidak terletak pada kerahasiaan
algoritma, melainkan sepenuhnya terletak pada kerahasiaan \textbf{kunci}. Asumsi
yang menyertainya tegas: penyerang mengetahui sepenuhnya algoritma kriptografi yang
digunakan, termasuk seluruh detail rancangannya.

Konsekuensi metodologis dari asumsi ini penting untuk disadari. Seorang perancang
tidak boleh mengandalkan kerumitan yang tidak terdokumentasi, nama variabel yang
membingungkan, atau trik yang hanya diketahui penulisnya sendiri. Yang harus
diandalkan adalah ketidakmampuan penyerang menyelesaikan persoalan matematis yang
mendasarinya. Inilah sebabnya keamanan kriptografi modern dirumuskan sebagai
pernyataan tentang kesulitan komputasi, bukan tentang ketidaktahuan penyerang.

Ada tiga cara lazim untuk mengelompokkan serangan, dan ketiganya akan dibahas pada
bab ini secara berurutan. Pengelompokan pertama mendasarkan diri pada
\textbf{informasi yang tersedia} bagi penyerang (Bagian~\ref{sec:jenis-serangan-informasi}),
kedua pada \textbf{tujuan atau hasil yang ingin dicapai} (Bagian~\ref{sec:tujuan-serangan}),
dan ketiga pada \textbf{keterlibatan penyerang} dalam komunikasi
(Bagian~\ref{sec:pasif-aktif}). Ketiga sudut pandang tersebut tidak saling
menggantikan; sebuah serangan nyata biasanya dapat dijelaskan secara bersamaan pada
ketiganya.
